\subsection{Минимальное межспутниковое расстояние}
\label{subsec:minimum_distance}

Помимо условий покрытия, к орбитальному построению целесообразно предъявлять ограничение на минимальное расстояние между любыми двумя КА. Обозначим через $r_{\min}^{\mathrm{Walker}}$ наименьшее евклидово расстояние между аппаратами группировки за один орбитальный период. Тогда условие безопасности можно записать в виде
\begin{equation}
    r_{\min}^{\mathrm{Walker}} \geq r_{\mathrm{req}},
\end{equation}
где $r_{\mathrm{req}}$ -- заданное допустимое минимальное межспутниковое
расстояние.

Для оценки $r_{\min}^{\mathrm{Walker}}$ предлагается использовать модель $J_2$: одинаковые $a$, $e=0$ и $i$ обеспечивают постоянство разностей ДВУ и аргументов широты внутри пары. Поэтому минимальное расстояние можно найти аналитически без численного распространения всей группировки.

\paragraph{Минимальное расстояние для пары КА}

Радиус-вектор КА на круговой орбите записывается как
\begin{equation}
\mathbf r(\Omega,u)=a
\begin{pmatrix}
\cos\Omega\cos u-\sin\Omega\cos i\sin u \\
\sin\Omega\cos u+\cos\Omega\cos i\sin u \\
\sin i\sin u
\end{pmatrix}.
\end{equation}
Рассмотрим два аппарата с разностью ДВУ $\Delta\Omega_{12}$ и постоянной разностью аргументов широты $\Delta u_{12}$. Благодаря осевой симметрии ДВУ первой орбиты можно принять равным нулю. В произвольный момент времени аргументы широты двух КА равны $u$ и
$u+\Delta u_{12}$, поэтому расстояние между ними определяется выражением
\begin{equation}
    d(u)=\left\|\mathbf r(0,u)-
    \mathbf r(\Delta\Omega_{12},u+\Delta u_{12})\right\|_2.
\end{equation}
Угол между нормалями к плоскостям пары
$i_{\mathrm{rel}}\in[0,\pi]$ определяется выражением
\begin{equation}
    \cos i_{\mathrm{rel}}=
    \cos^2 i+\sin^2 i\cos\Delta\Omega_{12}.
\end{equation}
Введём также вспомогательный угол
\begin{equation}
    \varkappa_{12}=\operatorname{atan2}\left(
    \cos\frac{\Delta\Omega_{12}}{2},
    \sin\frac{\Delta\Omega_{12}}{2}\cos i
    \right).
\end{equation}
Расстояние между КА равно модулю разности их радиус-векторов. Квадрат модуля разности двух векторов $\mathbf x$ и $\mathbf y$ равен $\|\mathbf x\|^2+\|\mathbf y\|^2-2\mathbf x\cdot\mathbf y$. Для рассматриваемой пары оба радиус-вектора имеют модуль $a$, поскольку рассматриваются круговые орбиты одинаковой высоты. В таком случае,
\begin{equation}
\label{eq:pair_d_u}
    d^2(u)=2a^2\left(1-
    \frac{\mathbf r(0,u)\cdot
    \mathbf r(\Delta\Omega_{12},u+\Delta u_{12})}{a^2}\right).
\end{equation}
Из \eqref{eq:pair_d_u} очевидно, что поиск минимального расстояния эквивалентен поиску максимума скалярного произведения радиус-векторов. В результате получаем
\begin{equation}
\label{eq:pair_rmin}
    r_{\min,12}=
    2a\left|\cos\left(\frac{\Delta u_{12}}{2}-\varkappa_{12}\right)\right|
    \cos\frac{i_{\mathrm{rel}}}{2}.
\end{equation}
Для плоскостей с разным ДВУ формула \eqref{eq:pair_rmin} получается из условия равенства нулю производной $d^2(u)$ по $u$, а положительный знак второй производной подтверждает, что стационарная точка \eqref{eq:pair_rmin} является минимумом. Для плоскостей с совпадающим ДВУ формула \eqref{eq:pair_rmin} сводится к длине хорды между двумя точками одной окружности:
\begin{equation}
    r_{\min,12}=2a
    \left|\sin\frac{\Delta u_{12}}{2}\right|.
\end{equation}

\paragraph{Минимум для Walker-построения}

Если принять ДВУ первой плоскости Walker-построения равным нулю, то для плоскости с индексом $p=1,\ldots,P$ и аппарата с индексом
$s=1,\ldots,S$ имеем
\begin{equation}
\begin{aligned}
    \Omega_p&=(p-1)\Delta\Omega,\\
    u_{p,s}(t)&=u(t)+(s-1)\frac{2\pi}{S}+(p-1)f.
\end{aligned}
\end{equation}
Следовательно, геометрия произвольной пары зависит не от самих индексов, а
только от их разностей:
\begin{equation}
\begin{aligned}
    \Delta p&=p_2-p_1,\\
    \Delta s&=s_2-s_1.
\end{aligned}
\end{equation}
\begin{equation}
\begin{aligned}
    \Delta\Omega_{12}&=\Delta p\Delta\Omega,\\
    \Delta u_{12}&=\Delta s\frac{2\pi}{S}+\Delta p f.
\end{aligned}
\end{equation}
Поэтому минимальное расстояние во всей группировке определяется как
\begin{equation}
\label{eq:walker_rmin}
    r_{\min}^{\mathrm{Walker}}=
    \min_{\substack{
        -(P-1)\leq\Delta p\leq P-1,\\
        -(S-1)\leq\Delta s\leq S-1,\\
        (\Delta p,\Delta s)\neq(0,0)}}
    r_{\min,12}\left(
        \Delta p\Delta\Omega,
        \Delta s\frac{2\pi}{S}+\Delta p f
    \right).
\end{equation}

Такой переход от пар аппаратов к разностям индексов исключает полный перебор всех пар. Вместо вычислительной сложности порядка $\mathcal{O}((PS)^2)$ необходимо проверить $(2P-1)(2S-1)-1$ комбинаций, то есть расчёт имеет сложность $\mathcal{O}(PS)$. При этом формула~\eqref{eq:walker_rmin} даёт минимум по периоду без дискретизации времени.

Для учёта ограничений безопасности в данном разделе было предложено оценивать $r_{\min}^{\mathrm{Walker}}$ в принятой приближённой модели движения. Такая оценка не заменяет более точного моделирования, однако её можно использовать при оптимизации, чтобы отдавать предпочтение построениям с большим номинальным разнесением КА, для которых можно ожидать большего минимального расстояния и в более точной модели.
